\section{Consequences for SAT and the limits of the improvement}
\label{sec:consequences}

\subsection{What changes after the existing reductions}

The source's arbitrary-density reduction uses a normalized wire budget
$s=O_k(n\log\log n)$ and other work of lower order than the displayed
worst-case exponent \cite{source}. Applying either new general
contraction algorithm retains
\[
 \poly_k(n,m)\,2^{O_k(n\log\log n/\sqrt{\log n})}.
\]
The polynomial-space branch retains polynomial space. The cached branch
uses subexponential space of the same scale as its time exponent.
For bounded clause density, the direct equality/clause network has
$O_{k,D}(n)$ wires and bounded degree, so the low-degree cached theorem
applies. Its order is still $2^{O_{k,D}(n/\sqrt{\log n})}$.

These consequences should not be expressed by replacing $s$ with $n$
inside a numerical coefficient. One must first pay the encoding's
wire expansion and normalize precision and tensor-description costs.
The density-sensitive result of the earlier packets can be retained
by taking the better applicable reduction \cite{density}.

\subsection{An explicit bounded-density wire count}
\label{sec:wirecount}

A simple count makes one translation entirely explicit.
Let $L$ be the number of literal occurrences after deleting duplicate
literals within clauses and tautologies. Let $d_v$ be the occurrence
degree of variable $v$. Make every clause vertex a left vertex. For
$d_v=1$, use a unary equality vertex on the right; for $d_v=2$, use
a binary equality vertex. For $d_v\ge3$, take an equality tree with
$d_v$ external ports and $d_v-2$ trivalent vertices, all provisionally
on the right. Its $d_v-3$ internal edges each receive a left binary
equality vertex. Clause-to-tree edges already have opposite colors,
so they require no subdivision.

The network is bipartite and has exactly
\begin{equation}
 M=L+2\sum_v(d_v-3)_+\le3L\le3km
 \label{eq:explicit-wires}
\end{equation}
wires and maximum degree $\max\{k,3\}$. Every consistent variable
value extends uniquely through its tree and the subdivisions.
Unused declared variables contribute the scalar multiplier
$2^{n_{\rm unused}}$. An empty clause is handled by returning zero;
an empty formula gives $2^n$.

Take $s=\max\{M,n+1,2\}$, so the modulus $2^{n+1}$ satisfies the
normalized precision premise. For fixed $k,D>0$ with $m\le Dn$,
put $B_{k,D}=\max\{3kD,1\}$. The low-degree result gives the explicit
conservative consequence
\[
 T(n,m)\le\poly_{k,D}(n,m)\,
 2^{(C_{\rm low}B_{k,D}+o(1))n/\sqrt{\log_2n}},
\]
and the corresponding space coefficient is
$C_{\rm low,space}B_{k,D}$. This is a direct wiring calculation,
not a claim that this encoding has the optimal wire constant.
It also does not provide a comparable numerical constant for the
source's general-density compression.

\subsection{Why the new theorems do not change logarithmic order}

At a call with $M$ wires and taper range $u$, the mode-entropy scale
is $M/\sqrt u$. The new summation arguments multiply that scale by
smaller constants. Their fixed group alphabets allow $u=\Theta(\log s)$
at the top of the recursion. They do not change either dependency.

This is why a substantial decrease in a coefficient can coexist with
an unchanged big-$O$ exponent class. If the objective is specifically
to replace the worst-case expression
$n\log\log n/\sqrt{\log n}$ by a little-$o$ expression, that objective
remains unresolved by this iteration.

\subsection{The stronger conditional proposals remain unproved}

The precision-refinement packet's target
$2^{O(n\log\log n/\log n)}$ requires a method summing the actual
structured mode expression with only $\exp(O(M/u))$ paid recursive
calls, preserving degree, precision, description size, and space.
Our adaptive separators still use $\exp(O(M/\sqrt u))$ calls.
The cached method reduces the exponent of its finite tensor summation,
but still spends $\exp(\Theta(M/\sqrt u))$ on its certified tables
in the worst envelope. Neither closes the missing interface.

The other conditional routes seek a low-degree representation after
compressing pair-mode lists or a uniformly linear-wire reduction.
Short binary indices alone do not preserve the bounded-degree
premise of the contraction theorem: a few tensors may carry a
large fraction of all the new wires. Treating their entries as
free recursive oracles would hide the main computation.

\section{Attempted routes to a smaller exponent}
\label{sec:attempts}

\subsection{Repeated mode compression and its budget obstruction}

One tempting approach is to compact the exact mode indices repeatedly.
There is a useful restricted warning. Suppose a fixed number $t$ of
successive, unfactored compact-mode constructions must obey taper
budgets of the form
\[
 u_1\log r_1\le A_1\log s,\qquad
 u_j\log r_j\le A_j\log r_{j-1}\quad(2\le j\le t),
\]
where $r_j>1$ and the $A_j$ are fixed constants. Multiplication gives
\begin{equation}
 \prod_{j=1}^t u_j
 \le \left(\prod_{j=1}^t A_j\right)
       \frac{\log s}{\log r_t}.
 \label{eq:telescope}
\end{equation}
Thus a modeled denominator $\sqrt{\prod_j u_j}$ is still at most
$O(\sqrt{\log s})$ when $r_t$ is bounded below by a constant greater
than one. Within this explicit multiplicative-budget model,
repeating the operation does not yield an extra power of $\log s$.

This is an algebraic implication of the stated premises, not a lower
bound for all compilers, all tensor algorithms, or SAT. The earlier
factorization and asymptotic-analysis packets explain how analogous
packing restrictions arise for their unfactored constructions
\cite{factorization,barriers}. A new factorization may invalidate
these premises. A growing number of stages also requires fresh
uniform control of the constants and is not covered by this
fixed-$t$ statement.

\subsection{Grouping parity tests before expanding incidences}

The general-density encoding has total retained incidence degree
$K=O(n)$ inside an incidence universe of size $N=n\,\polylog n$.
Its subset indexing costs
$O(K\log(1+N/K))=O(n\log\log n)$ wires.
We examined merging parity tests into fewer groups before this
expansion. Reducing the incidence universe is attractive, but a
single merged group can receive essentially all the allowed Taylor
or nilpotent degree. Its requested coefficient may then involve
all original variables. No uniformly shrinking coefficient oracle
was established in this review.

It would be incorrect to assume that oracle and count only the
smaller outer network. A successful proof must either exclude
such concentration, exploit a new algebraic representation, or
show a decreasing measure for the coefficient computation itself.
We retain this as a research question, not as a theorem.

\section{Further research questions}
\label{sec:further}

\begin{enumerate}
\item \textbf{Use more structure than pair entropy.}
The linear-flag theorem is sharp for balanced binary trees under
its Euclidean envelope, as proved in the preceding section.
Can the actual packet factors rule out its extremal entropy
configuration, or permit an unbalanced or non-tree summation?
This asks for a verifiable structural statement about generated
tensors, not a generic black-box improvement.

\item \textbf{Remove table construction as the cached bottleneck.}
The centered triangle estimate already puts the final triangle
below the largest binary block merge. Further improving only
the final multiplication will not lower the current local
certificate. A useful next result would reduce the $1+2$
merge score $916/25$ at nine groups, while continuing to charge
all singleton child entries.

\item \textbf{Construct the missing degree-preserving interface.}
A compact mode-index network needs efficient local factors with
controlled degree and immutable creation precision. Supplying
that interface could change the logarithmic order. Merely
showing that the index strings are short cannot.

\item \textbf{Reduce the incidence universe without an expensive oracle.}
Can a grouped parity lift compute concentrated coefficient
queries within a strictly decreasing size/precision measure?
A positive answer would address the extra $\log\log n$ factor
in the general-density reduction.

\item \textbf{Optimize a resource frontier with certified schedules.}
The present algorithms select clear fixed alphabets and rational
tapers. A finite dynamic program over valid alphabet changes,
tree choices, and memory caps could certify additional coefficient
improvements. Such optimization would refine the exponential
constant unless it introduces a new dependence on $u$.

\item \textbf{Formalize the finite new lemmas first.}
The coset-intersection lemma, the linear-flag averaging identity,
and the centered triangle estimate are compact mathematical
targets. Formalizing them would give reusable trusted
components before attempting the much larger analytic and
arithmetic integration.
\end{enumerate}
